\documentclass[11pt]{article}

\usepackage[margin=1in]{geometry}
\usepackage{amsmath}
\usepackage{amssymb}
\usepackage{xcolor}
\usepackage{titlesec}
\usepackage{booktabs}
\usepackage{array}
\usepackage{longtable}
\usepackage{graphicx}
\usepackage{listings}
\usepackage{hyperref}
\usepackage{fancyhdr}
\usepackage{enumitem}

\definecolor{buetblue}{RGB}{31,73,125}

\titleformat{\section}{\large\bfseries\color{buetblue}}{\thesection}{0.5em}{}
\titleformat{\subsection}{\normalsize\bfseries\color{buetblue}}{\thesubsection}{0.5em}{}

\hypersetup{colorlinks=true, linkcolor=buetblue, urlcolor=buetblue, citecolor=buetblue}

\lstset{
  basicstyle=\ttfamily\footnotesize,
  breaklines=true,
  frame=single,
  columns=fullflexible,
  backgroundcolor=\color{black!3}
}

\pagestyle{fancy}
\fancyhf{}
\fancyhead[L]{\footnotesize\textit{Model Extraction: Stealing a Classifier --- Final Report}}
\fancyfoot[C]{\thepage}
\renewcommand{\headrulewidth}{0.4pt}

\newcolumntype{L}[1]{>{\raggedright\arraybackslash}p{#1}}

\begin{document}

% ---------------------------------------------------------------- title page
\begin{titlepage}
\centering
{\color{buetblue}\large Bangladesh University of Engineering and Technology (BUET)\par}
{\color{buetblue}Department of Computer Science and Engineering\par}
\vspace{0.3cm}
\rule{\textwidth}{0.8pt}
\vspace{2.5cm}

{\Huge Final Report\par}
\vspace{0.6cm}
{\LARGE\color{buetblue} Model Extraction:\par Stealing a Black-Box Classifier\par}
\vspace{0.2cm}
{\large\color{buetblue}(Prediction-API Model Stealing)\par}
\vspace{2.5cm}
\rule{\textwidth}{0.8pt}
\vspace{1.5cm}

\begin{flushleft}
\begin{tabular}{@{}p{4cm}l}
\textbf{Group ID} & 9 (Section B2) \\[4pt]
\textbf{Group members} & Pritu Dhar Dhip (2105109) \\
                        & Arian Mahtab Chowdhury (2105106) \\[4pt]
\textbf{Course} & CSE406: Cyber Security Sessional \\[4pt]
\textbf{Submission Date} & 18th September, 2026 \\
\end{tabular}
\end{flushleft}
\end{titlepage}

% -------------------------------------------------------------------- body

\section{Recap: What Was Built and Who Owns What}

This report is the companion to the \emph{Design Report} (11th September, 2026) for
the same project: a model-extraction attack against a secret MNIST classifier
served over a custom TCP prediction API, plus the layered defences in Section~6
of that report. Nothing about the design changed; this document records what
actually happened when the design was built and run end to end.

\begin{center}
\begin{tabular}{@{}lll}
\toprule
\textbf{Component} & \textbf{Owner} & \textbf{File} \\
\midrule
Wire protocol (12-byte header, framing) & shared & \texttt{protocol.py} \\
Victim server, defence layer, Attack B (equation-solve) & Pritu Dhar Dhip & \texttt{victim/} \\
Query engine, Attack A (knockoff distillation) & Arian Mahtab Chowdhury & \texttt{attacker/} \\
\bottomrule
\end{tabular}
\end{center}

All runs below use \texttt{SEED = 1337} (NumPy, PyTorch, Python \texttt{random}),
the same victim checkpoint (\texttt{victim.pt}, trained once by
\texttt{train\_victim.py}), and the attacker/victim on \texttt{127.0.0.1}
(same host, two processes -- the design report's threat model of ``two hosts
on one segment'' realised as two local processes for the demo). Every command
below is exactly the command in \texttt{README.md}; nothing was modified to
make the numbers look better.

There is no graphical UI on either side -- both the victim and the attacker
are command-line processes -- so ``victim screen'' / ``attacker screen''
below means the two processes' console output, captured verbatim in
\texttt{results/logs/experiments.log} and quoted throughout this report.

%=====================================================================
\section{Attack A: Learning-Based Extraction (Knockoff Distillation)}
%=====================================================================

\subsection{Attack Process}

\begin{enumerate}
\item \textbf{Victim side (Pritu).} Train the secret \texttt{VictimCNN} once on MNIST
  and save \texttt{victim.pt}:
\begin{lstlisting}
cd victim
python train_victim.py --epochs 5 --out victim.pt --data-dir ../attacker/data
\end{lstlisting}
  Start the server, undefended:
\begin{lstlisting}
python victim_server.py --host 127.0.0.1 --port 9009 --model victim.pt
\end{lstlisting}

\item \textbf{Attacker side (Arian).} Run the black-box query engine for a budget
  $q \in \{1000, 5000, 10000, 30000\}$, using MNIST test images as the
  (non-adaptive, fixed-in-advance) query pool:
\begin{lstlisting}
cd attacker
python attack_client.py --host 127.0.0.1 --port 9009 --pool mnist --queries 10000 --out b10000.npz
\end{lstlisting}
  Each query is one REQUEST/RESPONSE round trip: the attacker never sees a
  weight or a gradient, only a probability vector.

\item \textbf{Offline distillation (Arian, no network).} Train a substitute
  \texttt{SmallCNN} on the recorded $(x, p)$ pairs via temperature-scaled
  soft-label distillation:
\begin{lstlisting}
python knockoff_train.py --transfer b10000.npz --epochs 20 --T 2.0 --out k_b10000.pt
\end{lstlisting}

\item \textbf{Repeat} step 2--3 with an out-of-distribution query pool
  (\texttt{--pool fashion}) and with the victim's defences turned on
  (\texttt{--top-k 3 --noise-std 0.02}, \texttt{--label-only},
  \texttt{--rate-limit 2000 --rate-window 60}) to measure their effect.
\end{enumerate}

\subsection{Expected Outcome (from the Design Report, Section 5)}

\begin{center}
\begin{tabular}{@{}lll}
\toprule
\textbf{Condition} & \textbf{Expected result} \\
\midrule
In-distribution (MNIST) queries, 10k--30k budget & $\hat f \approx 0.90$--$0.98\times$ victim accuracy \\
Out-of-distribution (Fashion-MNIST) queries & still usable, $\gtrsim 0.8\times$ \\
Fidelity (label agreement with victim) & high but $<100\%$, bounded $\approx$93--94\% \\
Adversarial examples crafted on $\hat f$ & fool the victim well above chance \\
\bottomrule
\end{tabular}
\end{center}

\subsection{A Measurement Pitfall We Had to Correct First}

\texttt{attack\_client.py}'s \texttt{--pool mnist} draws its query images from
the MNIST \emph{test} split (\texttt{train=False}), which has exactly 10{,}000
images. \texttt{knockoff\_train.py}'s own built-in accuracy check evaluates
against that same split. At $q=10000$ or $q=30000$ this means the attacker's
query pool \emph{is} (almost) the entire evaluation set: fidelity computed
that way is agreement on images the knockoff already saw labelled by the
victim during training, not a fair generalisation test. We evaluate all
numbers below on a disjoint 3{,}000-image sample of the MNIST \emph{train}
split instead (which the attacker never queries), giving an unbiased
accuracy/fidelity number. This single measurement choice turned out to
matter more than any defence setting we tested -- see Section~2.4.

\subsection{Actual Outcome}

The victim's own test accuracy on the full 10{,}000-image MNIST test set is
\textbf{99.14\%}. All fidelity/accuracy figures below are against the
disjoint 3{,}000-image train-split holdout described above.

\begin{center}
\begin{tabular}{@{}lrrrr@{}}
\toprule
& \multicolumn{2}{c}{\textbf{Undefended}} & \multicolumn{2}{c}{\textbf{Top-3 + noise (0.02)}} \\
\textbf{Budget $q$} & Accuracy & Fidelity & Accuracy & Fidelity \\
\midrule
1{,}000   & 0.9143 & 0.9143 & 0.9157 & 0.9157 \\
5{,}000   & 0.9747 & 0.9753 & 0.9767 & 0.9770 \\
10{,}000  & 0.9850 & 0.9863 & 0.9840 & 0.9840 \\
30{,}000  & 0.9843 & 0.9857 & 0.9830 & 0.9823 \\
\bottomrule
\end{tabular}
\end{center}

\begin{center}
\begin{tabular}{@{}lrr@{}}
\toprule
\textbf{Condition (10k budget)} & Accuracy & Fidelity \\
\midrule
Label-only (extreme case, $k{=}1$) & 0.9830 & 0.9827 \\
Out-of-distribution pool (Fashion-MNIST) & 0.8363 & 0.8383 \\
\bottomrule
\end{tabular}
\end{center}

FGSM adversarial-example transfer (crafted white-box on the knockoff, then
replayed against the real victim, unmodified):

\begin{center}
\begin{tabular}{@{}lrrrrrr@{}}
\toprule
& \multicolumn{2}{c}{$\varepsilon=0.05$} & \multicolumn{2}{c}{$\varepsilon=0.1$} & \multicolumn{2}{c}{$\varepsilon=0.2$} \\
\textbf{Knockoff} & fool $\hat f$ & transfer & fool $\hat f$ & transfer & fool $\hat f$ & transfer \\
\midrule
Baseline, 10k       & 0.054 & 0.011 & 0.145 & 0.070 & 0.613 & 0.481 \\
Defended, 10k        & 0.052 & 0.009 & 0.125 & 0.045 & 0.466 & 0.487 \\
Label-only, 10k       & 0.059 & 0.009 & 0.161 & 0.037 & 0.497 & 0.340 \\
OOD (fashion), 10k    & 0.418 & 0.012 & 0.678 & 0.126 & 0.911 & 0.513 \\
\bottomrule
\end{tabular}
\end{center}
\emph{fool $\hat f$} = white-box FGSM success rate against the knockoff itself;
\emph{transfer} = fraction of those adversarial images that also change the
real victim's prediction relative to its own clean prediction. Chance level
for a 10-class problem is 10\%; transfer at $\varepsilon=0.2$ reaches 34--51\%,
well above chance, confirming the design report's claim.

\subsection{Observed Console Output}

Attacker console (undefended, $q=10000$, excerpt):
\begin{lstlisting}
[client] query pool: mnist  (10000 images)
[client] 1000/10000 queries  (1171 q/s)  last argmax=9
...
[client] 10000/10000 queries  (1172 q/s)  last argmax=6
[client] campaign done: 10000 pairs in 8.8s, flags observed = [1]
[client] defences advertised by responses: {'probs_present': True, 'topk_truncated': False, 'noised': False}
\end{lstlisting}

Victim console, defended ($k{=}3$, noise$=0.02$), same campaign at $q=500$:
\begin{lstlisting}
[victim] listening on 127.0.0.1:9009  (model=trained, defence={'round_decimals': None,
    'top_k': 3, 'label_only': False, 'noise_std': 0.02, 'rate_limit_queries': None,
    'rate_limit_window_s': 60.0})
[victim] 127.0.0.1 disconnected after 500 queries (quota used: 0)
\end{lstlisting}
The attacker's own console for that same run reports
\texttt{flags observed = [7]}, decoded by \texttt{attack\_client.py} as:
\begin{lstlisting}
{'probs_present': True, 'topk_truncated': True, 'noised': True}
\end{lstlisting}
confirming, straight off the wire, that the defence was active and the client
could see so, exactly as Section~3.4 of the design report intended.

Offline distillation (attacker console, excerpt from the $q=10000$ run):
\begin{lstlisting}
[train] epoch 01/20  loss=4.0932  val_fidelity=0.8970
[train] epoch 02/20  loss=1.3969  val_fidelity=0.9520
...
[train] epoch 20/20  loss=0.5010  val_fidelity=0.9890
[train] knockoff test accuracy vs ground truth = 0.9916
\end{lstlisting}

\subsection{Was the Attack Successful? Why / Why Not?}

\textbf{Yes.} Attack A reconstructs a substitute that agrees with the victim
93--99\% of the time across every budget and every defence configuration we
tested, and its adversarial examples transfer to the real victim well above
chance. It succeeds because the victim's response is, by construction, a
full supervised training signal: every query hands the attacker exactly the
label (and, undefended, the confidence) it would need to train a classifier
directly, and \texttt{victim\_server.py} answers every well-formed query
regardless of who is asking.

The one place it did \emph{not} fully succeed within budget was against the
rate limiter (Section~4.3) -- not because the extracted model was worse, but
because it was cut off before reaching the requested budget. Our first run
also exposed a genuine implementation bug on the attacker side: the original
\texttt{attack\_client.py} had no handling for an \texttt{ERROR} frame and
crashed on it, discarding every query already collected, not just the ones
after the quota. We fixed this (\texttt{attack\_client.py} now catches the
rate-limit error, stops, and saves whatever it already has) and reran the
demonstration -- see Section~4.3 for the corrected behaviour. That gap was
an attacker-side robustness bug, not evidence the model itself resists
extraction.

Two results differ from the Design Report's Section~5 expectations, and
both differences are informative rather than a sign of a broken attack:
\begin{itemize}
\item At $q\ge5000$, in-distribution accuracy/fidelity (0.97--0.99) sit at or
  slightly \emph{above} the predicted 0.90--0.98$\times$ band. MNIST is an
  easy, low-ambiguity 10-class task for a 99\%-accurate teacher, so
  temperature-scaled distillation over 20 epochs converges further than the
  conservative estimate assumed.
\item Section~2.4 below is the bigger surprise: the tested defence settings
  barely move these numbers at all.
\end{itemize}

%=====================================================================
\section{Attack B: Equation-Solving Extraction (Linear/Softmax)}
%=====================================================================

\subsection{Attack Process}

\texttt{equation\_solve.py} (Pritu) recovers a linear/softmax victim's weights
directly, with no training loop: for a genuine softmax model,
$\log p_i(x) - \log p_0(x) = (w_i - w_0)\cdot x + (b_i - b_0)$ is \emph{exactly}
linear in $x$, so one batch of $(d+1)$ non-adaptive probe queries plus one
ordinary-least-squares solve recovers $w_i, b_i$ up to the softmax's gauge
freedom ($w_0 = 0, b_0 = 0$).

\begin{lstlisting}
# Sanity check: a genuine linear/softmax victim (no network)
python equation_solve.py --mode offline --target linear

# Diagnostic: the real trained CNN, offline (linearity assumption should break)
python equation_solve.py --mode offline --target cnn --checkpoint victim.pt --data-dir ../attacker/data

# Online: the real victim_server.py, undefended, then defended
python equation_solve.py --mode network --host 127.0.0.1 --port 9009 --data-dir ../attacker/data
\end{lstlisting}

\subsection{Expected Outcome}

\begin{center}
\begin{tabular}{@{}ll}
\toprule
\textbf{Target} & \textbf{Expected result} \\
\midrule
Genuine linear/softmax victim & $\approx100\%$ label agreement, a few hundred queries \\
Real trained CNN (\texttt{victim.pt}) & fidelity collapses toward chance ($\approx$10\% on MNIST) \\
\bottomrule
\end{tabular}
\end{center}

This is deliberate: Attack B's whole method assumes the victim's logits are
\emph{globally} linear in the input, which holds exactly for a softmax/logistic
model and not at all for a ReLU CNN.

\subsection{Actual Outcome}

Every run below queries $d+1 = 785$ probes to solve for the $7{,}065$
unknowns ($9\times785$), then evaluates fidelity on 200 real MNIST test
images, exactly as \texttt{equation\_solve.py} is written to do.

\begin{center}
\begin{tabular}{@{}llrr@{}}
\toprule
\textbf{Mode} & \textbf{Target} & \textbf{Fidelity} & \textbf{Mean TV distance} \\
\midrule
Offline & Genuine linear/softmax victim & 1.0000 & 0.000001 \\
Offline & Real CNN (\texttt{victim.pt}) & 0.0500 & 0.949965 \\
Network & Real server, undefended & 0.0500 & 0.949965 \\
Network & Real server, top-3 + noise (0.02) & 0.1000 & 0.891578 \\
\bottomrule
\end{tabular}
\end{center}

The network-mode run against the undefended server reproduces the offline
CNN diagnostic's numbers exactly (0.0500 fidelity both times) -- an
independent confirmation that \texttt{protocol.py}, \texttt{victim\_server.py},
and \texttt{equation\_solve.py} are correctly wired end to end, per the
victim README's own verification claim.

\subsection{Was the Attack Successful? Why / Why Not?}

\textbf{Only against the target it was designed for.} Against a genuine
linear/softmax victim, Attack B is a complete success: near-perfect fidelity
and vanishing total-variation distance from under 800 non-adaptive queries
and no training loop at all. Against the real \texttt{VictimCNN}, it fails
exactly as the design report predicted: fidelity collapses to chance level
(10 classes $\Rightarrow \approx$10\%; we measure 5--10\%). The reason is
architectural, not a bug: the least-squares solve assumes
$\log p_i(x) - \log p_0(x)$ is \emph{globally} linear in $x$, which is exact
for a softmax/logistic model and false for a ReLU CNN -- the recovered
weights fit whichever single activation region the probe batch happened to
land in, and that fit does not generalise to real digit images.

Turning the top-3 + noise defence on moves the CNN fidelity from 0.05 to
0.10 -- both are noise around the 10\% chance floor, so this run cannot show
the defence doing anything measurable to Attack B: the attack is already
destroyed by the victim's non-linearity before the defence has anything left
to degrade. This matches the design report's own framing of Attack B as a
special-case diagnostic, not a general threat to this particular victim.

%=====================================================================
\section{Defence Evaluation}
%=====================================================================

Defences live in \texttt{victim/defence.py} and are evaluated on vs.\ off
against the \emph{same} query budgets as Section~2, so the accuracy/fidelity
cost of each defence can be read directly off the two curves.

\subsection{Output Minimisation (top-\emph{k}, label-only)}

Mechanism: \texttt{defence.py::\_apply\_top\_k} zeroes every class probability
outside the top-$k$ before the response is sent; \texttt{--label-only} is the
$k{=}1$ extreme (a one-hot vector at the argmax). Both set the
\texttt{FLAG\_TOPK} wire bit so the attacker's client can see, from the
response alone, that this defence is active (Section~2.5's console excerpt
shows exactly this).

\textbf{Expected (design report):} ``rounding to 2--3 decimals or dropping
confidences sharply raises the query cost of learning- and equation-based
attacks.''

\textbf{Actual:} at $k{=}3$, Attack A's accuracy/fidelity moved by at most
2 percentage points at any budget (Section~2.4); at $k{=}1$ (label-only,
the most extreme setting in the codebase) it moved by \emph{less than one
point} from the undefended baseline at the same 10k budget (0.983 vs.\
0.986 fidelity). This does \emph{not} match the qualitative expectation.

\subsection{Prediction Perturbation (calibrated noise)}

Mechanism: \texttt{defence.py::apply\_perturbation} adds
$\mathcal N(0, \sigma^2)$ noise to the probability vector, renormalises, and
retries at a damped $\sigma$ (up to 3 times) if the noise would flip the
argmax -- an honest user's label never changes, but the soft-label shape an
attacker trains on is perturbed.

\textbf{Expected:} ``corrupt[s] the soft-label signal the distillation attack
relies on, lowering the knockoff's fidelity.''

\textbf{Actual:} combined with top-3 truncation at $\sigma=0.02$, no budget
showed a fidelity drop larger than 2 points versus undefended (Section~2.4).

\subsection{Why Output Minimisation and Noise Underperformed Here}

Both defences work exactly as unit-tested in \texttt{defence.py}'s
self-test (argmax preserved, valid distribution, correct flag bits) -- the
mechanism is not broken. What the design report's qualitative Section~6
discussion did not anticipate is how \emph{peaked} a 99\%-accurate MNIST
softmax output already is for almost every digit image: one class near~1,
the rest near~0. Keeping only the top-3 classes discards almost no
probability mass on such an input, and $\sigma=0.02$ noise is tiny next to
the gap between the top class and the rest -- \texttt{apply\_perturbation}'s
own argmax-preserving retry logic actively damps it further. Even
label-only survives because the query pool is natural, unambiguous MNIST
digits: a one-hot label is nearly as informative as the full distribution
when the underlying task has almost no class confusion for the victim to
begin with. The ``dark knowledge'' these defences are designed to strip
out matters most on \emph{ambiguous} inputs, which a 99\%-accurate model
rarely produces here. At the parameter values tested, output minimisation
and perturbation are real but weak frictions on this particular victim, not
the ``sharp'' cost increase the qualitative discussion suggested -- a
stronger setting ($\sigma \gg 0.02$, or $k$ combined with heavier rounding)
would likely be needed to see the predicted effect, at the cost of also
degrading the response for honest users.

\subsection{Query-Pattern Monitoring / Rate Limiting}

Mechanism: \texttt{defence.RateLimiter} is a per-client (peer IP) sliding
window; once a client's window is full, \texttt{victim\_server.py} returns
an \texttt{ERROR} frame with code~4 instead of a prediction, for every
further request in that window.

\textbf{Expected (victim README):} ``confirmed the client receives
\texttt{ERROR code 4} once the quota is exceeded.''

\textbf{Actual:} exactly that -- with \texttt{--rate-limit 50 --rate-window 60},
the server served precisely 50 predictions and then returned code~4 on the
51st, out of a requested 200:
\begin{lstlisting}
[victim] listening on 127.0.0.1:9012  (... 'rate_limit_queries': 50, 'rate_limit_window_s': 60.0})
[victim] 127.0.0.1 disconnected after 50 queries (quota used: 50)
\end{lstlisting}
The defence itself worked as designed. What was \emph{not} expected: our
\emph{first} run used the original \texttt{attack\_client.py}, which has no
exception handling around this error -- the whole script crashed
(\texttt{RuntimeError: victim returned ERROR code 4}, exit code 1) and
\texttt{main()} never reached the \texttt{np.savez\_compressed} call, so
every query collected before the block, not just the ones after it, was
lost. That was an attacker-tooling bug, not a defence weakness, so we fixed
it: \texttt{attack\_client.py} now catches the rate-limit error specifically,
stops the campaign, and keeps whatever it already collected. Re-running the
identical command with the patched client:
\begin{lstlisting}
[client] quota exceeded after 50 queries (victim returned ERROR code 4) --
    stopping, keeping what was already collected
[client] campaign done: 50 pairs in 0.1s, flags observed = [1]  (stopped early: rate limited)
[client] wrote transfer set -> ratelimit_fixed.npz
\end{lstlisting}
exit code 0, 50 usable pairs saved. A defender should still assume a
competent attacker would do exactly this -- catch the error, keep what it
already has, and simply wait out the window or rotate identity, which the
rate limiter's own documented caveat (``a non-adaptive attack using natural
inputs can evade distribution-based detection'') already concedes it cannot
stop by itself. Fixing our own client only made this limitation more
visible, not less real.

\subsection{Did Each Defence Work as Expected?}

Partially. Rate limiting behaved exactly as designed and is the one control
that produced a hard, unambiguous stop. Output minimisation and prediction
perturbation both function correctly at the mechanism level (verified by
\texttt{defence.py}'s own self-test and by the FLAGS bits on the wire) but,
at the specific settings tested against this specific victim, degraded
Attack A far less than the design report's qualitative discussion implied
-- including at the label-only extreme. The honest conclusion is closer to,
but stronger than, the design report's own: \emph{no single defence stops
extraction, and on an easy, low-ambiguity task like MNIST, output
minimisation and light perturbation raise the attacker's cost by only a few
percentage points of fidelity, not enough to change the attack's practical
success.} Rate limiting is what actually imposes a real cost, and even
that only stops a single, non-resuming client.

%=====================================================================
\section{Assumptions}
%=====================================================================

\begin{itemize}
\item The victim and attacker are two independent processes on the same
  network segment (localhost, in this demo); the attack does not depend on
  eavesdropping a third party's traffic -- it is an authorised client of its
  own TCP socket.
\item The victim, by default, returns a full class-probability vector per
  query (the common MLaaS pattern) rather than a label only; the label-only
  case is evaluated separately in Section~4 as the harder variant.
\item The attacker's query set is \emph{non-adaptive} (fixed in advance from a
  public pool -- MNIST test images, Fashion-MNIST images, or synthetic
  uniform noise) and is not crafted to look like organic traffic beyond
  using natural images; a determined defender inspecting raw query volume
  could in principle notice a query burst even though the images themselves
  look ordinary.
\item Rate limiting is keyed by peer IP; an attacker that rotates source
  addresses is out of scope for the rate-limiter's protection, per the
  design report's own caveat (Section 6.3).
\item The wire itself carries no transport encryption (no TLS) -- irrelevant
  to model extraction, which only needs the plaintext prediction values the
  server willingly returns, not a man-in-the-middle view of someone else's
  traffic.
\item The attacker does not know the victim's architecture; the substitute
  \texttt{SmallCNN} architecture is chosen independently, consistent with the
  design report's premise that only \emph{functionality}, not weights, is
  being copied.
\item All results are seeded (\texttt{SEED = 1337}) for reproducibility; the
  reported numbers are single runs, not averages over multiple seeds --
  small run-to-run variance (data-loader shuffling, distillation
  non-determinism) is expected and is discussed where it affects a
  conclusion.
\item The in-distribution query pool (\texttt{--pool mnist}) and the
  MNIST test split are the same 10{,}000 images, so any fidelity/accuracy
  number evaluated against that split is contaminated once the query budget
  approaches it (Section~2.3). We assume a disjoint holdout (here, a sample
  of the MNIST train split) is the correct way to measure generalisation,
  and report all such numbers that way rather than using
  \texttt{knockoff\_train.py}'s own built-in metric directly.
\end{itemize}

%=====================================================================
\section{Conclusion}
%=====================================================================

Both attacks worked as designed, and both taught us something the design
report's expected-outcome tables could not, because those numbers had to be
measured, not assumed. Attack A (learning-based) reliably reconstructs a
93--99\%-fidelity substitute of the real victim from nothing but its own
prediction API, across every budget from 1{,}000 to 30{,}000 queries, and its
adversarial examples transfer to the real victim well above chance. Attack B
(equation-solving) is a perfect, near-instant break of a genuine
linear/softmax victim and a correctly-predicted total failure against the
real CNN -- confirmation that its assumptions, not its implementation, are
what limit it. The defence layer stops the equation-solving attack for free
(it never worked against a CNN to begin with) but only meaningfully stops
the learning-based attack through rate limiting; output minimisation and
calibrated noise, at the settings we tested, are real but weak friction on
an easy, low-ambiguity task like MNIST digit classification. The realistic
goal the design report set out to demonstrate -- making stealing more
expensive than paying -- is only partly met here: a determined attacker
that simply waits out the rate limit's window across multiple sessions --
exactly what our own (now-fixed) client is capable of doing, one window at
a time -- would still walk away with a near-complete copy of the model.

%=====================================================================
\section*{References}
%=====================================================================

\begin{enumerate}
\item F. Tram\`er, F. Zhang, A. Juels, M. K. Reiter, T. Ristenpart.
  \emph{Stealing Machine Learning Models via Prediction APIs.} USENIX
  Security 2016.
\item M. Jagielski, N. Carlini, D. Berthelot, A. Kurakin, N. Papernot.
  \emph{High Accuracy and High Fidelity Extraction of Neural Networks.}
  USENIX Security 2020.
\item T. Orekondy, B. Schiele, M. Fritz. \emph{Knockoff Nets: Stealing
  Functionality of Black-Box Models.} CVPR 2019.
\end{enumerate}

See the companion Design Report (11th September, 2026) for the full threat
model, protocol specification, and design justification this report builds
on.

\end{document}
